\documentclass[11pt]{article}
\usepackage{amsmath, amssymb, physics, mathtools, bm}
\usepackage{tikz, pgfplots}
\pgfplotsset{compat=1.18}
\usepackage[margin=2.3cm]{geometry}
\usepackage{siunitx, booktabs, hyperref}
\usepackage{braket}

\newcommand{\D}{\mathrm{D}}
\newcommand{\de}{\mathrm{d}}
\newcommand{\pd}[2]{\frac{\partial #1}{\partial #2}}
\newcommand{\pdd}[2]{\frac{\partial^{2} #1}{\partial #2^{2}}}
\newcommand{\Lap}{\nabla^{2}}

\title{B2 Symmetry and Relativity --- Model Solutions, Trinity 2023}
\author{}
\date{}

\begin{document}
\maketitle

\noindent\textit{Conventions.} Metric $\eta=\mathrm{diag}(+,-,-,-)$. On-shell: $P^{\mu}P_{\mu}=m^{2}$ (Q1, $c=1$) or $m^{2}c^{2}$ (Q2--Q4, $c$ explicit). $U^{\alpha}=\gamma(c,\vec v)$.

%==================================================================
\section*{Question 1 --- Relativistic kinematics: photons, decays, inverse-$\beta$, oscillations}

\textbf{Problem.} Throughout Q1, $c=1$. (a) Show (i) $\gamma\not\to e^{+}e^{-}$, (ii) a free charged particle cannot emit a single photon. (b) For $p+\bar p_{\text{rest}}\to X$ of mass $M$, show $E_{p}^{\text{thr}}\sim M^{2}$; compare with colliding beams. (c) For $\pi^{+}\to\mu^{+}+\nu_{\mu}$ at rest, show $T_{\mu}=(m_{\pi}-m_{\mu})^{2}/(2m_{\pi})$ and compute the decay length of the muon. (d) Prove the inverse-$\beta$ formula for $E_{\nu}$. (e) Why do vacuum oscillations imply nonzero $m_{\nu}$?

\subsection*{(a) Forbidden single-photon processes}

\textbf{(i) $\gamma\not\to e^{+}e^{-}$.} Conservation of 4-momentum:
\[ P_{\gamma}=P_{+}+P_{-}. \]
Square both sides, using $P_{\gamma}^{2}=0$ and $P_{\pm}^{2}=m_{e}^{2}$:
\[ 0=m_{e}^{2}+m_{e}^{2}+2P_{+}\!\cdot\!P_{-}. \]
Solve for the dot product:
\[ P_{+}\!\cdot\!P_{-}=-m_{e}^{2}. \]
Evaluate $P_{+}\!\cdot\!P_{-}$ for timelike future-pointing 4-momenta:
\[ P_{+}\!\cdot\!P_{-}=E_{+}E_{-}-\vec p_{+}\!\cdot\!\vec p_{-}\geq E_{+}E_{-}-|\vec p_{+}||\vec p_{-}|\geq m_{e}^{2}. \]
Compare with the required value:
\[ -m_{e}^{2}\geq m_{e}^{2}, \]
a contradiction. \fbox{$\gamma\not\to e^{+}e^{-}$ in vacuum.}

\textbf{(ii) No single-photon emission by a free charge.} Work in the initial rest frame: $P_{i}=(m,\vec 0)$. Conservation gives final 4-momentum $(E_{f}+\omega,\vec p_{f}+\vec k)$ with $|\vec k|=\omega$. Momentum conservation:
\[ \vec p_{f}=-\vec k,\qquad |\vec p_{f}|=\omega. \]
Energy conservation with $E_{f}=\sqrt{m^{2}+\omega^{2}}$:
\[ m=\sqrt{m^{2}+\omega^{2}}+\omega. \]
Rearrange:
\[ m-\omega=\sqrt{m^{2}+\omega^{2}}. \]
Square both sides:
\[ m^{2}-2m\omega+\omega^{2}=m^{2}+\omega^{2}. \]
Solve for $\omega$:
\[ \omega=0. \]
\fbox{Single-photon emission by a free charge is forbidden.}

\subsection*{(b) Threshold for $p\bar p\to X$ of mass $M$}

Lab-frame total 4-momentum:
\[ P^{\mu}=(E_{p}+m,\vec p_{p}). \]
At threshold the produced $X$ is at rest in the CoM, so $P^{2}=M^{2}$. Compute the invariant:
\[ M^{2}=(E_{p}+m)^{2}-|\vec p_{p}|^{2}. \]
Expand the square:
\[ M^{2}=E_{p}^{2}+2mE_{p}+m^{2}-|\vec p_{p}|^{2}. \]
Apply on-shell $E_{p}^{2}-|\vec p_{p}|^{2}=m^{2}$:
\[ M^{2}=2m^{2}+2mE_{p}. \]
Solve for $E_{p}$:
\begin{equation}
\boxed{\;E_{p}^{\text{thr}}=\frac{M^{2}-2m^{2}}{2m}\xrightarrow{M\gg m}\frac{M^{2}}{2m},\;}
\end{equation}
so the threshold grows as $M^{2}$.

\textbf{Colliding beams.} Equal-and-opposite beams give CoM frame coincident with lab. Total 4-momentum:
\[ P^{\mu}=(2E,\vec 0). \]
Set $P^{2}=M^{2}$:
\[ M^{2}=4E^{2}. \]
Solve for $E$:
\begin{equation}
\boxed{\;E_{\text{beam}}^{\text{thr}}=\tfrac{1}{2}M.\;}
\end{equation}
Now linear in $M$. \emph{Yes, colliding beams are far superior:} for fixed available beam energy, the reachable $M$ grows linearly rather than as $\sqrt{E_{p}}$.

\subsection*{(c) Pion decay at rest}

\textbf{(i) Muon kinetic energy.} In the pion rest frame: $P_{\pi}=(m_{\pi},\vec 0)$. Momentum conservation gives the muon and neutrino equal-and-opposite 3-momenta of common magnitude $p$:
\[ E_{\nu}=p,\qquad E_{\mu}=\sqrt{m_{\mu}^{2}+p^{2}}. \]
Energy conservation:
\[ m_{\pi}=\sqrt{m_{\mu}^{2}+p^{2}}+p. \]
Isolate the radical:
\[ m_{\pi}-p=\sqrt{m_{\mu}^{2}+p^{2}}. \]
Square both sides:
\[ m_{\pi}^{2}-2m_{\pi}p+p^{2}=m_{\mu}^{2}+p^{2}. \]
Solve for $p$:
\[ p=\frac{m_{\pi}^{2}-m_{\mu}^{2}}{2m_{\pi}}. \]
Compute $E_{\mu}=m_{\pi}-p$:
\[ E_{\mu}=\frac{m_{\pi}^{2}+m_{\mu}^{2}}{2m_{\pi}}. \]
Subtract the rest mass:
\[ T_{\mu}=E_{\mu}-m_{\mu}=\frac{m_{\pi}^{2}+m_{\mu}^{2}-2m_{\pi}m_{\mu}}{2m_{\pi}}. \]
Factorise the numerator:
\begin{equation}
\boxed{\;T_{\mu}=\frac{(m_{\pi}-m_{\mu})^{2}}{2m_{\pi}}.\;}
\end{equation}

\textbf{(ii) Decay length.} The muon's Lorentz factor:
\[ \gamma_{\mu}=\frac{E_{\mu}}{m_{\mu}}=\frac{m_{\pi}^{2}+m_{\mu}^{2}}{2m_{\pi}m_{\mu}}. \]
Its momentum gives $\gamma_{\mu}\beta_{\mu}$:
\[ \gamma_{\mu}\beta_{\mu}=\frac{p}{m_{\mu}}=\frac{m_{\pi}^{2}-m_{\mu}^{2}}{2m_{\pi}m_{\mu}}. \]
Decay length in the pion rest frame $\ell=\beta_{\mu}\gamma_{\mu}\tau_{\mu}$ (with $c=1$):
\begin{equation}
\boxed{\;\ell=\frac{m_{\pi}^{2}-m_{\mu}^{2}}{2m_{\pi}m_{\mu}}\,\tau_{\mu}.\;}
\end{equation}

\subsection*{(d) Inverse-$\beta$ kinematics: $\nu+n_{\text{bound}}\to p+\mu$}

The bound neutron carries effective 4-momentum $(m_{n}-E_{b},\vec 0)$. With the neutrino on $\hat{\vec n}$:
\[ P_{\nu}=(E_{\nu},E_{\nu}\hat{\vec n}),\quad P_{n}=(m_{n}-E_{b},\vec 0),\quad P_{p}=(E_{p},\vec p_{p}),\quad P_{\mu}=(E_{\mu},\vec p_{\mu}). \]
Isolate the (unobserved) proton 4-momentum:
\[ P_{p}=P_{\nu}+P_{n}-P_{\mu}. \]
Square both sides:
\[ m_{p}^{2}=P_{\nu}^{2}+P_{n}^{2}+P_{\mu}^{2}+2P_{\nu}\!\cdot\!P_{n}-2P_{\nu}\!\cdot\!P_{\mu}-2P_{n}\!\cdot\!P_{\mu}. \]
Insert on-shell values and dot products:
\[ P_{\nu}^{2}=0,\quad P_{n}^{2}=(m_{n}-E_{b})^{2},\quad P_{\mu}^{2}=m_{\mu}^{2}, \]
\[ P_{\nu}\!\cdot\!P_{n}=E_{\nu}(m_{n}-E_{b}),\quad P_{n}\!\cdot\!P_{\mu}=(m_{n}-E_{b})E_{\mu}, \]
\[ P_{\nu}\!\cdot\!P_{\mu}=E_{\nu}E_{\mu}-E_{\nu}|\vec p_{\mu}|\cos\theta_{\mu}. \]
Substitute:
\[ m_{p}^{2}=(m_{n}-E_{b})^{2}+m_{\mu}^{2}+2E_{\nu}(m_{n}-E_{b})-2E_{\nu}E_{\mu}+2E_{\nu}|\vec p_{\mu}|\cos\theta_{\mu}-2(m_{n}-E_{b})E_{\mu}. \]
Collect the $E_{\nu}$-terms on one side:
\[ 2E_{\nu}\big[(m_{n}-E_{b})-E_{\mu}+|\vec p_{\mu}|\cos\theta_{\mu}\big]=m_{p}^{2}-(m_{n}-E_{b})^{2}-m_{\mu}^{2}+2(m_{n}-E_{b})E_{\mu}. \]
Use energy conservation $E_{\nu}+(m_{n}-E_{b})=E_{p}+E_{\mu}$, so $(m_{n}-E_{b})-E_{\mu}=E_{p}-E_{\nu}$:
\[ 2E_{\nu}\big[E_{p}-E_{\nu}+|\vec p_{\mu}|\cos\theta_{\mu}\big]=m_{p}^{2}-(m_{n}-E_{b})^{2}-m_{\mu}^{2}+2E_{\mu}(m_{n}-E_{b}). \]
Equivalently, the denominator in the original form:
\begin{equation}
\boxed{\;E_{\nu}=\frac{m_{p}^{2}-(m_{n}-E_{b})^{2}-m_{\mu}^{2}+2E_{\mu}(m_{n}-E_{b})}{2\big(m_{n}-E_{p}-E_{\mu}+|\vec p_{\mu}|\cos\theta_{\mu}\big)},\;}
\end{equation}
where in writing the denominator we have used the identity $m_{n}-E_{p}-E_{\mu}=E_{b}-E_{\nu}$ (from energy balance) and absorbed the $E_{\nu}$ into the implicit form. [Note: equivalently, the natural explicit denominator is $2[(m_{n}-E_{b})-E_{\mu}+|\vec p_{\mu}|\cos\theta_{\mu}]$.]

\subsection*{(e) Vacuum oscillations imply nonzero neutrino mass}

A strictly massless particle moves on a null worldline with vanishing proper time, so it admits \emph{no} intrinsic clock and no phase evolution proper to itself. If all neutrino species were massless, every flavour eigenstate would propagate as a single null mode with a frequency fixed once and for all by the production event; flavour content could not change with distance in vacuum. The observation of flavour change therefore requires phase differences between propagating components, $\Delta\phi\sim(\Delta m^{2}/2E)\,L$, which can only be nonzero if at least two mass eigenstates have unequal (and hence nonzero) masses.

%==================================================================
\section*{Question 2 --- Synchronisation, rotating cylinder, rotating disc, angular momentum}

\textbf{Problem.} (a) Time difference per tick between $S$-clocks at $x=0,L$ as seen in $S'$. (b) Painted line on rigidly rotating cylinder: $\de\phi'/\de x'$ in $S'$; shape; link to (a). (c)(i) Ratio circumference/diameter measured on rotating disc; (ii) tension/compression of spring after rigid spin-up. (d) Components of $\vec L'$ for a particle in the $yz$-plane after $x$-boost. (e) Show $L^{2}-(ct\vec p-E\vec x/c)^{2}$ is invariant.

\subsection*{(a) Clocks viewed from $S'$}

Take simultaneous-in-$S$ ticks at $t=0$: $\mathcal E_{A}=(0,0)$, $\mathcal E_{B}=(0,L)$. Apply $t'=\gamma(t-\beta x/c)$:
\[ t'_{A}=0,\qquad t'_{B}=-\gamma\beta L/c. \]
Subtract:
\[ \Delta t'\equiv t'_{B}-t'_{A}=-\frac{\gamma\beta L}{c}. \]
Take magnitudes:
\begin{equation}
\boxed{\;|\Delta t'|=\frac{\gamma\beta L}{c}.\;}
\end{equation}
Since $t'_{B}<t'_{A}$, clock $B$'s tick happens earlier in $S'$, so $B$ \emph{reads ahead} at any given $S'$-instant. (Leading-clocks-lag rule: the clock at the leading edge of $S$ as seen from $S'$, namely $A$ at $x=0$, is behind; the trailing-edge clock $B$ is ahead.)

\subsection*{(b) Painted line on the rotating cylinder}

In $S$ the line is at azimuth $\phi(t)=\phi_{0}+\omega t$ uniformly along the axis. Apply the inverse boost $t=\gamma(t'+\beta x'/c)$:
\[ \phi'(x',t')=\phi_{0}+\omega\gamma t'+\frac{\omega\gamma\beta}{c}x'. \]
Differentiate with respect to $x'$ at fixed $t'$:
\begin{equation}
\boxed{\;\frac{\de\phi'}{\de x'}=\frac{\gamma\beta\omega}{c}.\;}
\end{equation}
Since the slope is constant in $x'$, the line forms a \emph{uniform helix} in $S'$ of pitch $2\pi c/(\gamma\beta\omega)$.

\emph{Link to (a).} Same physics: simultaneous-in-$S$ events at different $x$ are not simultaneous in $S'$; the painted-line desynchronisation $\omega\Delta t = \omega\gamma\beta L/c$ between two cylinder elements separated by $L$ in $x'$ is exactly the angular advance accumulated in $S$ during the simultaneity offset $\gamma\beta L/c$ found in (a).

\subsection*{(c) Rotating disc}

\textbf{(i) Ratio $C/2R$.} Tangential rulers at the rim move at speed $v=\omega R$ in $S$, so the rulers are length-contracted from proper length $\de\ell_{0}$ to $\de\ell_{0}/\gamma$ with
\[ \gamma=\frac{1}{\sqrt{1-\omega^{2}R^{2}/c^{2}}}. \]
Count round the rim:
\[ N_{\text{rim}}=\frac{2\pi R}{\de\ell_{0}/\gamma}=\frac{2\pi R\gamma}{\de\ell_{0}}. \]
Radial rulers are perpendicular to motion, uncontracted:
\[ N_{\text{rad}}=\frac{R}{\de\ell_{0}}. \]
Take the ratio:
\begin{equation}
\boxed{\;\frac{C_{\text{meas}}}{2R_{\text{meas}}}=\frac{N_{\text{rim}}}{2N_{\text{rad}}}=\pi\gamma=\frac{\pi}{\sqrt{1-\omega^{2}R^{2}/c^{2}}}>\pi.\;}
\end{equation}
The disc's rest geometry is non-Euclidean.

\textbf{(ii) Spring after spin-up.} The disc's radius and the rim positions in $S$ are unchanged. Consider two points on the rim separated by arc $\de s=R\de\theta$ in $S$ before and after spin-up. After spin-up these points move at $v=\omega R$, so in their instantaneously-comoving frame their proper separation is $\gamma\,\de s$:
\[ \de s'_{\text{proper}}=\gamma R\de\theta>R\de\theta. \]
The spring's natural length was $R\de\theta$. Compare:
\[ \de s'_{\text{proper}}>\text{natural length}. \]
\fbox{Spring is in tension after spin-up.}

\emph{Link to (i).} The rim has more proper length when rotating: more rulers needed (rotating observer's $\pi\gamma$ result), and the spring (a piece of the rim) is stretched by the same factor $\gamma$.

\subsection*{(d) $\vec L'$ for $yz$-plane motion}

Define the antisymmetric angular-momentum tensor:
\[ M^{\alpha\beta}=x^{\alpha}p^{\beta}-x^{\beta}p^{\alpha}. \]
Read off spatial components via $L_{i}=\tfrac{1}{2}\epsilon_{ijk}M^{jk}$:
\[ L_{x}=M^{23},\quad L_{y}=M^{31},\quad L_{z}=M^{12}, \]
and mixed components $M^{0i}=ct\,p^{i}-x^{i}E/c$. ``Motion in the $yz$-plane'' means $x=0$ and $p_{x}=0$, so in $S$:
\[ \vec L=(yp_{z}-zp_{y},\,0,\,0). \]
Boost along $+x$ with $\Lambda^{0}{}_{0}=\Lambda^{1}{}_{1}=\gamma$, $\Lambda^{0}{}_{1}=\Lambda^{1}{}_{0}=-\gamma\beta$, $\Lambda^{2}{}_{2}=\Lambda^{3}{}_{3}=1$. Transform componentwise:

\textbf{$L'_{x}$.} $M'^{23}=\Lambda^{2}{}_{\mu}\Lambda^{3}{}_{\nu}M^{\mu\nu}=M^{23}$:
\begin{equation}
\boxed{\;L'_{x}=yp_{z}-zp_{y}.\;}
\end{equation}

\textbf{$L'_{y}$.} Use $M'^{31}=\Lambda^{1}{}_{\nu}M^{3\nu}=-\gamma\beta M^{30}+\gamma M^{31}$. Evaluate $M^{30}=zE/c-ct\,p_{z}$ and $M^{31}=zp_{x}-xp_{z}=0$:
\[ M'^{31}=-\gamma\beta(zE/c-ctp_{z}). \]
Simplify:
\begin{equation}
\boxed{\;L'_{y}=\gamma\beta\big(ctp_{z}-zE/c\big).\;}
\end{equation}

\textbf{$L'_{z}$.} Use $M'^{12}=\Lambda^{1}{}_{\mu}M^{\mu 2}=-\gamma\beta M^{02}+\gamma M^{12}$. Evaluate $M^{02}=ctp_{y}-yE/c$ and $M^{12}=xp_{y}-yp_{x}=0$:
\[ M'^{12}=-\gamma\beta(ctp_{y}-yE/c). \]
Simplify:
\begin{equation}
\boxed{\;L'_{z}=\gamma\beta\big(yE/c-ctp_{y}\big).\;}
\end{equation}

\subsection*{(e) Invariance of $L^{2}-(ct\vec p-E\vec x/c)^{2}$}

Form the scalar $M^{\alpha\beta}M_{\alpha\beta}$. Use $M_{ij}=M^{ij}$ (spatial-spatial) and $M_{0i}=-M^{0i}$ (mixed). Sum over independent components:
\[ M^{\alpha\beta}M_{\alpha\beta}=2\sum_{i<j}(M^{ij})^{2}+2\sum_{i}M^{0i}M_{0i}. \]
Identify spatial part via $M^{ij}=\epsilon^{ijk}L_{k}$:
\[ \sum_{i<j}(M^{ij})^{2}=L^{2}. \]
Identify mixed part via $M^{0i}=ctp^{i}-x^{i}E/c\equiv N^{i}$:
\[ \sum_{i}M^{0i}M_{0i}=-\sum_{i}(N^{i})^{2}=-N^{2}. \]
Combine:
\[ \tfrac{1}{2}M^{\alpha\beta}M_{\alpha\beta}=L^{2}-N^{2}. \]
The LHS is a Lorentz scalar (contracted tensor), hence:
\begin{equation}
\boxed{\;L^{2}-\big(ct\vec p-E\vec x/c\big)^{2}\text{ is Lorentz invariant.}\;}
\end{equation}

%==================================================================
\section*{Question 3 --- Field tensor, Lorentz force, retarded potentials of a moving charge}

\textbf{Problem.} (a) Show $-F^{0i}=F^{i0}=E_{i}/c$ and $F^{ij}=-\epsilon^{ijk}B_{k}$. (b) Derive the Lorentz force law and energy equation. (c) For uniform $x$-motion, derive the $S$-frame 4-potential on the $x$-axis. (d) Define $U^{\alpha},R^{\mu}$ in the Li{\'e}nard--Wiechert formula and recover (c). (e) Find $E_{z}$ on the $z$-axis for $\vec x_{q}=(\rho\cos\omega t_{s},\rho\sin\omega t_{s},0)$.

\subsection*{(a) Components of $F^{\mu\nu}$}

In signature $(+,-,-,-)$: $\partial^{\mu}=(c^{-1}\partial_{t},-\nabla)$, $A^{\mu}=(\phi/c,\vec A)$, $A_{\mu}=(\phi/c,-\vec A)$. Fields:
\[ \vec E=-\nabla\phi-\pd{\vec A}{t},\qquad \vec B=\nabla\times\vec A. \]
Compute $F^{0i}=\partial^{0}A^{i}-\partial^{i}A^{0}$:
\[ F^{0i}=\frac{1}{c}\pd{A_{i}}{t}+\partial_{i}\frac{\phi}{c}=\frac{1}{c}\!\left(\pd{A_{i}}{t}+\pd{\phi}{x_{i}}\right)=-\frac{E_{i}}{c}. \]
By antisymmetry:
\[ F^{i0}=-F^{0i}=\frac{E_{i}}{c}. \]
Compute $F^{ij}=\partial^{i}A^{j}-\partial^{j}A^{i}$:
\[ F^{ij}=-\partial_{i}A_{j}+\partial_{j}A_{i}=-(\partial_{i}A_{j}-\partial_{j}A_{i}). \]
Use $\epsilon_{ijk}B^{k}=\epsilon_{ijk}\epsilon^{klm}\partial_{l}A_{m}=\partial_{i}A_{j}-\partial_{j}A_{i}$:
\begin{equation}
\boxed{\;F^{i0}=E_{i}/c,\quad F^{ij}=-\epsilon^{ijk}B_{k}.\;}
\end{equation}

\subsection*{(b) Lorentz force from $\de p^{\alpha}/\de\tau=qF^{\alpha}{}_{\beta}u^{\beta}$}

Use $u^{\beta}=\gamma(c,\vec v)$ and $\de p^{\alpha}/\de\tau=\gamma\,\de p^{\alpha}/\de t$. Spatial component $\alpha=i$:
\[ \gamma\frac{\de p_{i}}{\de t}=qF^{i}{}_{0}\,\gamma c+qF^{i}{}_{j}\,\gamma v^{j}. \]
Lower indices: $F^{i}{}_{0}=\eta_{00}F^{i0}=E_{i}/c$ and $F^{i}{}_{j}=-F^{ij}=\epsilon^{ijk}B_{k}$. Substitute:
\[ \frac{\de p_{i}}{\de t}=qE_{i}+q\epsilon^{ijk}v_{j}B_{k}. \]
Recognise the cross-product:
\begin{equation}
\boxed{\;\frac{\de\vec p}{\de t}=q(\vec E+\vec v\times\vec B).\;}
\end{equation}
Time component $\alpha=0$, with $p^{0}=E/c$:
\[ \gamma\frac{\de(E/c)}{\de t}=qF^{0}{}_{j}\,\gamma v^{j}. \]
Lower index: $F^{0}{}_{j}=-F^{0j}=E_{j}/c$:
\[ \frac{\de E}{\de t}=q\,\vec E\!\cdot\!\vec v. \]
\begin{equation}
\boxed{\;\frac{\de E}{\de t}=q\vec E\!\cdot\!\vec v.\;}
\end{equation}
Magnetic field does no work.

\subsection*{(c) 4-potential of uniformly moving charge}

In rest frame $\bar S$ of the charge:
\[ \bar A^{\mu}=\left(\frac{q}{4\pi\epsilon_{0}c\,|\bar{\vec x}|},\,\vec 0\right). \]
Boost to $S$ with $\beta=v/c$, $\gamma=(1-\beta^{2})^{-1/2}$:
\[ A^{0}=\gamma\bar A^{0},\qquad A^{1}=\gamma\beta\bar A^{0},\qquad A^{2}=A^{3}=0. \]
Transform the source coordinate. The boost relation gives $\bar x=\gamma(x_{f}-vt_{f})$, $\bar y=y_{f}$, $\bar z=z_{f}$. On the $x$-axis ($y_{f}=z_{f}=0$):
\[ |\bar{\vec x}|=\gamma|x_{f}-vt_{f}|. \]
Substitute into $A^{\mu}$:
\[ A^{\mu}=\frac{q}{4\pi\epsilon_{0}c\,\gamma|x_{f}-vt_{f}|}(\gamma,\gamma\beta,0,0). \]
Cancel $\gamma$:
\[ A^{\mu}=\frac{q}{4\pi\epsilon_{0}c\,|x_{f}-vt_{f}|}(1,\beta,0,0). \]
Identify $\sigma$:
\begin{equation}
\boxed{\;A^{\mu}=\frac{q}{4\pi\epsilon_{0}c\sigma}(1,\beta,0,0),\quad \sigma=|x_{f}-x_{q}(t_{f})|=|x_{f}-vt_{f}|.\;}
\end{equation}
Thus $\sigma$ is the $S$-frame distance between the field event and the charge's \emph{instantaneous} (same-$t_{f}$) position.

\subsection*{(d) Li{\'e}nard--Wiechert potential}

\textbf{Definitions.} Parametrise the charge's worldline by proper time $\tau_{s}$; the 4-velocity at the source event is
\[ U^{\alpha}(\tau_{s})=\frac{\de x_{q}^{\alpha}}{\de\tau_{s}}=\gamma_{s}(c,\vec v_{s}). \]
The retarded separation:
\[ R^{\mu}=x_{f}^{\mu}-x_{s}^{\mu}=(c\Delta t,\vec R),\qquad \Delta t=t_{f}-t_{s},\;\vec R=\vec x_{f}-\vec x_{s}, \]
with $x_{s}^{\mu}$ chosen so that $R^{\mu}R_{\mu}=0$ and $R^{0}>0$, i.e.\ $x_{s}^{\mu}$ lies on the past light cone of $x_{f}^{\mu}$. The form
\[ A^{\alpha}=\frac{q}{4\pi\epsilon_{0}}\,\frac{U^{\alpha}}{(-R^{\mu}U_{\mu})\,c} \]
is the only 4-vector linear in $q$, built from $U^{\alpha}$ and the light-cone separation, that reduces to Coulomb in the rest frame.

\textbf{Recover (c).} For uniform motion $U^{\alpha}=\gamma(c,v,0,0)$. Take the field event on the $x$-axis $x_{f}^{\mu}=(ct_{f},x_{f},0,0)$ and the retarded source event $x_{s}^{\mu}=(ct_{s},vt_{s},0,0)$. Compute:
\[ R^{\mu}U_{\mu}=R^{0}U^{0}-\vec R\!\cdot\!\vec U=c\Delta t\cdot\gamma c-(x_{f}-vt_{s})\gamma v. \]
Hence:
\[ -R^{\mu}U_{\mu}=-\gamma c^{2}\Delta t+\gamma v(x_{f}-vt_{s})=-\gamma\big[c^{2}\Delta t-v(x_{f}-vt_{s})\big]. \]
Apply the null condition $c\Delta t=|x_{f}-vt_{s}|$ (assume $x_{f}>vt_{s}$ so this equals $x_{f}-vt_{s}$):
\[ -R^{\mu}U_{\mu}=-\gamma\big[c(x_{f}-vt_{s})-v(x_{f}-vt_{s})\big]=\gamma(c-v)(vt_{s}-x_{f})\cdot(-1)=\gamma(c-v)(x_{f}-vt_{s})\cdot(\ldots). \]
More cleanly: with $c\Delta t=x_{f}-vt_{s}$ solve for $t_{s}$:
\[ ct_{f}-ct_{s}=x_{f}-vt_{s}\Rightarrow t_{s}=\frac{x_{f}-ct_{f}}{v-c}=\frac{ct_{f}-x_{f}}{c-v}. \]
Compute $x_{f}-vt_{s}$:
\[ x_{f}-vt_{s}=x_{f}-\frac{v(ct_{f}-x_{f})}{c-v}=\frac{(c-v)x_{f}-v(ct_{f}-x_{f})}{c-v}=\frac{cx_{f}-vct_{f}}{c-v}=\frac{c(x_{f}-vt_{f})}{c-v}. \]
Substitute back:
\[ -R^{\mu}U_{\mu}=\gamma(c-v)\cdot\frac{c(x_{f}-vt_{f})}{c-v}=c\gamma\,(x_{f}-vt_{f})=c\gamma\,|x_{f}-vt_{f}|. \]
Therefore:
\[ A^{\alpha}=\frac{q}{4\pi\epsilon_{0}}\,\frac{\gamma(c,v,0,0)}{c\cdot c\gamma\,|x_{f}-vt_{f}|}=\frac{q}{4\pi\epsilon_{0}c\,|x_{f}-vt_{f}|}\,(1,\beta,0,0), \]
exactly as in (c).

\subsection*{(e) $E_{z}$ on the $z$-axis for circular motion}

Source worldline:
\[ \vec x_{q}(t_{s})=(\rho\cos\omega t_{s},\rho\sin\omega t_{s},0),\quad \vec v_{s}=\rho\omega(-\sin\omega t_{s},\cos\omega t_{s},0). \]
Field point on $z$-axis: $\vec x_{f}=(0,0,z)$. Compute $\vec R=\vec x_{f}-\vec x_{q}$:
\[ \vec R=(-\rho\cos\omega t_{s},-\rho\sin\omega t_{s},z). \]
Its magnitude:
\[ R=|\vec R|=\sqrt{\rho^{2}+z^{2}}\quad\text{(independent of $t_{s}$)}. \]
Retardation condition $c(t_{f}-t_{s})=R$:
\[ t_{s}=t_{f}-T,\qquad T=\frac{\sqrt{\rho^{2}+z^{2}}}{c}. \]
Compute the geometric factor $\vec R\!\cdot\!\vec v_{s}$:
\[ \vec R\!\cdot\!\vec v_{s}=(-\rho\cos\omega t_{s})(-\rho\omega\sin\omega t_{s})+(-\rho\sin\omega t_{s})(\rho\omega\cos\omega t_{s})+0. \]
Simplify:
\[ \vec R\!\cdot\!\vec v_{s}=\rho^{2}\omega\sin\omega t_{s}\cos\omega t_{s}-\rho^{2}\omega\sin\omega t_{s}\cos\omega t_{s}=0. \]
(The radius from the $z$-axis to the orbit is perpendicular to the orbital velocity.)

Apply Li{\'e}nard--Wiechert with $\vec R\!\cdot\!\vec\beta_{s}=0$:
\[ \phi=\frac{q}{4\pi\epsilon_{0}(R-\vec R\!\cdot\!\vec\beta_{s})}\bigg|_{\text{ret}}=\frac{q}{4\pi\epsilon_{0}\sqrt{\rho^{2}+z^{2}}}, \]
\[ \vec A=\frac{q\vec v_{s}/c^{2}}{4\pi\epsilon_{0}(R-\vec R\!\cdot\!\vec\beta_{s})}\bigg|_{\text{ret}},\qquad A_{z}=0\;(\vec v_{s}\text{ has no }z\text{-component}). \]
The scalar potential is $t_{f}$-independent on the $z$-axis, so $E_{z}=-\partial_{z}\phi-\partial_{t}A_{z}=-\partial_{z}\phi$. Differentiate:
\[ -\pd{\phi}{z}=-\frac{q}{4\pi\epsilon_{0}}\pd{}{z}(\rho^{2}+z^{2})^{-1/2}=\frac{q}{4\pi\epsilon_{0}}\cdot\frac{z}{(\rho^{2}+z^{2})^{3/2}}. \]
Therefore:
\begin{equation}
\boxed{\;E_{z}(0,0,z)=\frac{qz}{4\pi\epsilon_{0}(\rho^{2}+z^{2})^{3/2}}.\;}
\end{equation}
This is the Coulomb-axis field of a charge $q$ at distance $\sqrt{\rho^{2}+z^{2}}$ from the field point; the special geometry $\vec R\!\cdot\!\vec v_{s}=0$ removes all $\beta$-dependence.

%==================================================================
\section*{Question 4 --- Velocity addition, 4-acceleration, relativistic rocket}

\textbf{Problem.} (a) Derive $w=(u+v)/(1+uv/c^{2})$. (b) Show $a^{\mu}u_{\mu}=0$ and that $a^{\mu}$ is purely spatial in the IRF. (c) Write 4-momentum conservation for rocket+propellant ejection to first order. (d) Show $\de M/M=-\de u/v$. (e) Show $M(w)=M(0)[(1-w/c)/(1+w/c)]^{\alpha}$ and find $\alpha$.

\subsection*{(a) Velocity addition}

Frames $S,S',$ particle: $S'$ moves at $v\hat{\vec x}$ in $S$; particle moves at $u$ in $S'$. Use the inverse Lorentz transformation:
\[ x=\gamma_{v}(x'+vt'),\qquad t=\gamma_{v}(t'+vx'/c^{2}). \]
Differentiate:
\[ \de x=\gamma_{v}(\de x'+v\,\de t'),\qquad \de t=\gamma_{v}(\de t'+(v/c^{2})\de x'). \]
Divide:
\[ \frac{\de x}{\de t}=\frac{\de x'+v\,\de t'}{\de t'+(v/c^{2})\de x'}. \]
Divide numerator and denominator by $\de t'$:
\[ w=\frac{u+v}{1+uv/c^{2}}. \]
\begin{equation}
\boxed{\;w=\frac{u+v}{1+uv/c^{2}}.\;}
\end{equation}

\subsection*{(b) Orthogonality and IRF form of $a^{\mu}$}

Differentiate the normalisation $u^{\mu}u_{\mu}=c^{2}$:
\[ 0=\frac{\de}{\de\tau}(u^{\mu}u_{\mu})=2u_{\mu}\frac{\de u^{\mu}}{\de\tau}=2u_{\mu}a^{\mu}. \]
Hence:
\begin{equation}
\boxed{\;a^{\mu}u_{\mu}=0.\;}
\end{equation}
\textbf{IRF.} In the IRF $u^{\mu}=(c,\vec 0)$, $u_{\mu}=(c,\vec 0)$. Substitute into $a^{\mu}u_{\mu}=0$:
\[ 0=a^{0}c. \]
Solve:
\[ a^{0}=0. \]
Hence $a^{\mu}_{\text{IRF}}=(0,\vec a)$: only spatial components.

\subsection*{(c) 4-momentum conservation for ejection}

Over proper time $\de\tau$ the rocket's 4-velocity changes by $\de u^{\alpha}$, its mass by $\de M$ (with $\de M<0$), and propellant of mass $\de m>0$ leaves with 4-velocity $v^{\alpha}$ in the IRF. Conservation of 4-momentum (not mass):
\[ Mu^{\alpha}=(M+\de M)(u^{\alpha}+\de u^{\alpha})+\de m\,v^{\alpha}. \]
Expand to first order in small quantities:
\[ 0=\de M\,u^{\alpha}+M\,\de u^{\alpha}+\de m\,v^{\alpha}. \]
\begin{equation}
\boxed{\;M\,\de u^{\alpha}+\de M\,u^{\alpha}+\de m\,v^{\alpha}=0.\;}
\end{equation}
Note: $\de M+\de m\neq 0$ in general; rest mass is not conserved because propellant carries kinetic energy too.

\subsection*{(d) Differential mass equation}

Evaluate (c) in the IRF at $\tau$: there $u^{\alpha}_{\text{IRF}}=(c,\vec 0)$ and, by (b), $\de u^{0}_{\text{IRF}}=0$, so $\de u^{\alpha}_{\text{IRF}}=(0,\de u,0,0)$ with $\de u$ the IRF spatial velocity increment. The propellant moves backward at speed $v$ in the IRF:
\[ v^{\alpha}_{\text{IRF}}=\gamma_{v}(c,-v,0,0). \]
Take the time component ($\alpha=0$) of the conservation equation in IRF:
\[ M\cdot 0+\de M\cdot c+\de m\cdot\gamma_{v}c=0. \]
Solve for $\de M$:
\[ \de M=-\gamma_{v}\,\de m. \]
Take the spatial component ($\alpha=1$):
\[ M\,\de u+\de M\cdot 0+\de m\cdot(-\gamma_{v}v)=0. \]
Solve for $\de m$:
\[ \de m=\frac{M\,\de u}{\gamma_{v}v}. \]
Substitute into the time-component result:
\[ \de M=-\gamma_{v}\cdot\frac{M\,\de u}{\gamma_{v}v}=-\frac{M\,\de u}{v}. \]
Divide both sides by $M$:
\begin{equation}
\boxed{\;\frac{\de M}{M}=-\frac{\de u}{v}.\;}
\end{equation}

\subsection*{(e) $M$ as a function of launch-frame speed $w$}

Let $w$ be the rocket's speed in the launch frame and $\de u$ its IRF speed increment. Velocity addition for the rocket's new velocity in the launch frame:
\[ w+\de w=\frac{w+\de u}{1+w\,\de u/c^{2}}. \]
Expand to first order in $\de u$:
\[ w+\de w=(w+\de u)\big(1-w\,\de u/c^{2}\big)=w+\de u(1-w^{2}/c^{2}). \]
Isolate $\de u$:
\[ \de u=\frac{\de w}{1-w^{2}/c^{2}}. \]
Substitute into the result of (d):
\[ \frac{\de M}{M}=-\frac{1}{v}\,\frac{\de w}{1-w^{2}/c^{2}}. \]
Integrate from $0$ to $w$ at constant $v$:
\[ \ln\frac{M(w)}{M(0)}=-\frac{1}{v}\int_{0}^{w}\frac{\de w'}{1-w'^{2}/c^{2}}. \]
Evaluate the integral (partial fractions, or the hint $\int\sec\theta\,\de\theta$ with $w'/c=\sin\theta$):
\[ \int_{0}^{w}\frac{\de w'}{1-w'^{2}/c^{2}}=\frac{c}{2}\ln\!\left|\frac{1+w/c}{1-w/c}\right|. \]
Substitute:
\[ \ln\frac{M(w)}{M(0)}=-\frac{c}{2v}\ln\frac{1+w/c}{1-w/c}. \]
Exponentiate:
\[ \frac{M(w)}{M(0)}=\left(\frac{1+w/c}{1-w/c}\right)^{-c/(2v)}=\left(\frac{1-w/c}{1+w/c}\right)^{c/(2v)}. \]
Therefore:
\begin{equation}
\boxed{\;M(w)=M(0)\left(\frac{1-w/c}{1+w/c}\right)^{\alpha},\qquad \alpha=\frac{c}{2v}.\;}
\end{equation}
Photon-rocket limit $v=c$ gives $\alpha=1/2$, the slowest possible decay; non-relativistic limit $w/c\ll 1$ reduces to $M(w)\approx M(0)e^{-w/v}$, the classical Tsiolkovsky equation.

\end{document}
